\subsection{\texorpdfstring{\(C_{l}\) 型系统 (\(l \geqslant 2\))}{C\_l 型系统 (l >= 2)}}

\begin{enumerate}
\def\labelenumi{(\Roman{enumi})}
\item
  \(C_{l}\) 型根系的存在已在第 5 号中得到证明，因为我们已经看到，\(B_{l}\) 型系统的逆系统是 \(C_{l}\) 型。因此，一个 \(C_{l}\) 型根系可由在 \(R^{l}\) 中取向量 \(\pm2\varepsilon_{i}\ (1\leqslant i\leqslant l)\) 和 \(\pm\varepsilon_{i}\pm\varepsilon_{j}\ (1\leqslant i<j\leqslant l)\) 得到。根的数目为 \(2l^{2}\)。
\item
  对第 5 号中所考察系统的基应用映射 \(\alpha\mapsto\frac{2\alpha}{(\alpha|\alpha)}\) 所得的像，可以作为 R 的一个基。我们得到：
\end{enumerate}

\[
\alpha_ {1} = \varepsilon_ {1} - \varepsilon_ {2}, \alpha_ {2} = \varepsilon_ {2} - \varepsilon_ {3}, \dots , \alpha_ {l - 1} = \varepsilon_ {l - 1} - \varepsilon_ {l}, \alpha_ {l} = 2 \varepsilon_ {l}.
\]

正根是 \(2\varepsilon_{i}\) 以及 \(\varepsilon_{i} \pm \varepsilon_{j} (i < j)\)。

\begin{enumerate}
\def\labelenumi{(\Roman{enumi})}
\setcounter{enumi}{2}
\item
  Coxeter 数与逆系统的相同：\(h = 2l\)。
\item
  令 \(\tilde{\alpha}=2\varepsilon_{1}=2\alpha_{1}+2\alpha_{2}+\cdots+2\alpha_{l-1}+\alpha_{l}\)，这是一个根。相对于 \((\alpha_{i})\)，它的坐标之和为 2l-1=h-1。因此，\(\tilde{\alpha}\) 是最高根。我们有 \((\tilde{\alpha}|\alpha_{i})=0\)（当 \(i\neq l\) 时），\((\tilde{\alpha}|\alpha_{l})=2\)。从而，扩展 Dynkin 图为
\end{enumerate}

\[
\alpha_ {1} \quad \alpha_ {2} \quad \dots \quad \alpha_ {l - 2} \quad \alpha_ {l - 1} \quad \alpha_ {l}
\]

\begin{enumerate}
\def\labelenumi{(\Alph{enumi})}
\setcounter{enumi}{21}
\tightlist
\item
  我们已经确定了 \(R^{-}\)，它是 \(B_{l}\) 型的。由 § 1 第 12 号的公式 (19) 及第 5 号 (V)，\(2\varepsilon_{i}\) 在 \(\Phi_{R}\) 下的长度平方为
\end{enumerate}

\[
((l + 1) (4 l - 2)) ^ {+ 1} ((4 l - 2) ^ {- 1}) ^ {- 1} = (l + 1) ^ {- 1};
\]

从而，\(\Phi_{\mathrm{R}}(x,y) = (x|y) / 4(l + 1)\)。

我们有 \(\gamma(R) = \gamma(R') = (l + 1)(4l - 2)\)。

\begin{enumerate}
\def\labelenumi{(\Roman{enumi})}
\setcounter{enumi}{5}
\tightlist
\item
  基本权容易求得：
\end{enumerate}

\[
\begin{array}{l} \omega_ {i} = \varepsilon_ {1} + \varepsilon_ {2} + \dots + \varepsilon_ {i} \\ = \alpha_ {1} + 2 \alpha_ {2} + \dots + (i - 1) \alpha_ {i - 1} + i (\alpha_ {i} + \alpha_ {i + 1} + \dots + \frac {1}{2} \alpha_ {l}) (i \leqslant l). \end{array}
\]

\begin{enumerate}
\def\labelenumi{(\Roman{enumi})}
\setcounter{enumi}{6}
\tightlist
\item
  正根之和为
\end{enumerate}

\[
\begin{array}{r l} 2 \rho & = 2 l \varepsilon_ {1} + (2 l - 2) \varepsilon_ {2} + \dots + 4 \varepsilon_ {l - 1} + 2 \varepsilon_ {l} \\ & = 2 l \alpha_ {1} + 2 (2 l - 1) \alpha_ {3} + \dots + i (2 l - i + 1) \alpha_ {i} + \dots \\ & \qquad \dots + (l - 1) (l + 2) \alpha_ {l - 1} + \frac {1}{2} l (l + 1) \alpha_ {l}. \end{array}
\]

\begin{enumerate}
\def\labelenumi{(\Roman{enumi})}
\setcounter{enumi}{7}
\item
  根据第 4 号和第 5 号 (VIII)，\(Q(R) = L_{1}\)，\(P(R) = L_{0}\)；\(P(R)/Q(R)\) 同构于 Z/2Z，连接指标为 2。
\item
  以及 (X) 这些数据只取决于 W(R)，因此与 \(B_{l}\) 型的相同。
\item
  与第 5 号中相同的论证表明 \(A(R) = W(R)\) 且 \(w_0 = -1\)。
\item
  \(\mathrm{P}(\mathrm{R}^{-})/\mathrm{Q}(\mathrm{R}^{-})\) 的唯一非单位元定义了扩展 Dynkin 图的唯一非平凡自同构：它把对应于 \(\alpha_{j}\) 和 \(\alpha_{l-j}\)（\(0 \leqslant j \leqslant l\)）的顶点互换。
\end{enumerate}

